%% GENERATED by scripts/paper/build_prose.py from scripts/paper/prose.yaml (the prose) and the analysis
%% register (the options blocks) -- edit the prose there, or here once you stop regenerating.

% ======== f-personal-style ========

\textcolor{red}{I have made a first pass. This is a very nice start. There are a handful of numbers to fill in, and one or two graphs to change. Make a note of what you once had in here - just because I've removed it doesn't mean I never want it, but the crucial stuff is now here. I think later on (we'll get to this), I may want to integrate some of the numbers and some of the quotes that you put in here before but I've since removed}

\subsection{Each operator has a personal style}\label{sec:res:personal-style}

Operators differ a great deal in how they use each assistant, and the differences appear personal. We describe each session by ten habits (how often a strategy card is chosen, how often a tactical suggestion is requested, how long a decision takes, etc). From those habits alone, a participant's second session resembles their own first (median profile $r = .48$) and essentially nobody else's ($r = -.01$): their own session is in the closest three out of the 33 people far more often than chance (41\% vs 9\%), even though the second session is a different scenario. Who the operator is explains approximately six times more of strategy card choice than the mission does. The styles do not fall into types, and nothing measured beforehand (attitude to AI, views on human authority, experience) predicts them.

% --- f-personal-style#habits (pass 5: numbers only) ---

\phantomsection\label{sec:res:personal-style:habits}
To compare styles we describe every 8-minute session by ten habits. Figure~\ref{fig:res:personal-style:habits} shows each habit's spread over all 66 sessions. Most habits vary enormously between sessions (strategy card choice runs from 0\% to 100\% of a session's decisions); two are nearly universal (most sessions keep every tactical suggestion unchanged, and most never edit a strategy card). The right-hand column shows the correlation (Spearman's rank) between each habit's session 1 and session 2; i.e.\ how much a person's habit in one run carries over to the next.

\begin{figure}[H]
\centering
\includegraphics[width=\linewidth,height=0.50\textheight,keepaspectratio]{figs/results/f-personal-style__habits__main.png}
\caption{\textbf{The habits describing each user and how they vary.} These ten habits accurately describe a user and many carry over to that user's following scenario.}
\label{fig:res:personal-style:habits}
\end{figure}

% --- f-personal-style#fingerprint-ranks (pass 5: for the paper) ---

\phantomsection\label{sec:res:personal-style:fingerprint-ranks}
When asked to describe their approach, 15 of 33 named a consistent strategy of their own. In Figure~\ref{fig:res:personal-style:fingerprint-ranks:all}, each person's ten-habit profile from session 1 is compared with all 33 session-2 profiles, and the 33 are ranked from most to least alike (and the same the other way round, session 2 among the session 1s) by Euclidean distance over the 10 dimensions. If habits were not personal, a person's own other session would be equally likely to land anywhere in that ranking (dashed line). Instead it is the single closest match 12 times out of 66 (2 expected by chance) and in the closest three 27 times (6 expected), even though the second session is a different scenario ($p < .001$).

\begin{figure}[H]
\centering
\includegraphics[width=\linewidth,height=0.50\textheight,keepaspectratio]{figs/results/f-personal-style__fingerprint-ranks__all.png}
\caption{\textbf{A person's own other session is usually among their closest matches.} Where each person's own session ranks among all 33 by ten-habit profile; dashed line: where it would fall if habits were not personal.}
\label{fig:res:personal-style:fingerprint-ranks}
\label{fig:res:personal-style:fingerprint-ranks:all}
\end{figure}

% --- f-personal-style#who-vs-what (pass 5: for the paper) ---

\phantomsection\label{sec:res:personal-style:who-vs-what}
When inspecting every instance of a user choosing or not choosing a strategy card, knowing \textit{who} is deciding predicts whether a card is chosen about six times better than knowing everything about the mission and when it arrived. We built three logistic models with access to different priors about the decision. Knowing only \textit{who} made the choice, the model classified the right one 75\% of the time (a coin: 50\%; knowing only the mission: 63\%). A person's session-1 rate on its own still does it 71\% of the time for their session-2 decisions. Figure~\ref{fig:res:personal-style:who-vs-what:r2} shows this as a set of $R^2$ values.

\begin{figure}[H]
\centering
\includegraphics[width=\linewidth,height=0.50\textheight,keepaspectratio]{figs/results/f-personal-style__who-vs-what__r2.png}
\caption{\textbf{Who decides explains far more of strategy card choice than what is being decided.} McFadden pseudo-$R^2$ of models predicting whether a strategy card is chosen from the mission and its timing, from the person, and from both.}
\label{fig:res:personal-style:who-vs-what}
\label{fig:res:personal-style:who-vs-what:r2}
\end{figure}

% --- f-personal-style#axis-scatter (pass 5: for the paper) ---

\phantomsection\label{sec:res:personal-style:axis-scatter}
Participants' strategic and tactical reliance correlated only slightly. Figure~\ref{fig:res:personal-style:axis-scatter:between} shows each participant's proportional use of each assistant. The correlation between use of the two assistants is only $\rho = 0.37$.

\begin{figure}[H]
\centering
\includegraphics[width=\linewidth,height=0.50\textheight,keepaspectratio]{figs/results/f-personal-style__axis-scatter__between.png}
\caption{\textbf{Use of one assistant says little about use of the other.} Each dot is a person: the share of their strategic decisions where a strategy card was chosen, against the share of their tactical plans where a suggestion was requested.}
\label{fig:res:personal-style:axis-scatter}
\label{fig:res:personal-style:axis-scatter:between}
\end{figure}

% --- f-personal-style#not-types (pass 5: numbers only) ---

\phantomsection\label{sec:res:personal-style:not-types}
The styles are individual rather than types; we were not able to statistically isolate the users into groups based on similar behaviours. Additionally, nothing measured before the study predicts a participant's style. Figure~\ref{fig:res:personal-style:not-disposition} shows that from 40 tests of attitude to AI (AIAS-4), verification propensity, views on human authority and prior experience against five behaviours, three reach $p < .05$ and none survives Holm correction. Although it should be noted that the variability in AIAS-4 was limited, this still suggests no statistically signifiant relationship (which we could find) between people's philosophy towards AI and their usage pattern in this work.

\begin{figure}[H]
\centering
\includegraphics[width=\linewidth,height=0.50\textheight,keepaspectratio]{figs/results/f-personal-style__not-disposition__main.png}
\caption{\textbf{Nothing measured beforehand predicts style.} Correlations between the pre-study scales and five behaviours; none survives correction for the 40 tests.}
\label{fig:res:personal-style:not-disposition}
\end{figure}

These findings suggest that the interaction style of a human user varies widely, and cannot easily be predicted by their attitudes or prior experiences. This suggests that the human-machine team formed in this type of problem depends heavily on which human is operating it, and as soon as the human role is handed over to a new operator, the behaviour of the swarm changes substantively.
\textcolor{red}{TODO in the discussion re this: Call for the importance of training and consistency, and possibly the use of agents to steer us towards this.}


\clearpage

% ======== f-size-decides ========

\subsection{Automation Preference to Simpler Tasks}\label{sec:res:size-decides}
A key question in this scenario is which tasks the human choses to delegate. We found that the bigger the job, the more likely an operator was to hand it to an assistant, on both tiers. Figure~\ref{fig:res:size-decides:both} shows the proportion of missions that the user delegated to the AI, grouped by the size in terms of drone requirement. As the mission size increases, so does the propensity to defer to the assistant. This is more pronounced at the tactical level than the strategic. Similar results (not displayed graphically here) are seen when comparing the number of different colours of drone required for a task, where missions including 3 different colours are automated X\% vs Y\% for those with only 2). This appeared to be an explicit aspect of operator philosophy, for example: ``I think if it's just one colour, I'm gonna do it, like, assign it myself, but if it's multiple colours, I think I'm gonna let, um, the AI choose for me, for the, the strategic side'' (P29), ``When they only needed a single colour of drone I tended to do it manually'' (P34).

\begin{figure}[H]
\centering
\includegraphics[width=\linewidth,height=0.50\textheight,keepaspectratio]{figs/results/f-size-decides__both__main.png}
\caption{\textbf{The bigger the mission, the more it is handed to an assistant, on both tiers.} Share of decisions where a strategy card was chosen (blue) or a tactical suggestion requested (orange), by the drones the mission needs; whiskers are 95\% intervals.}
\label{fig:res:size-decides:both}
\end{figure}

% --- f-size-decides#models (pass 5: for the paper) ---

\phantomsection\label{sec:res:size-decides:models}
The effect holds when each person is compared only with themselves, for every way of measuring size and on both tiers. Figure~\ref{fig:res:size-decides:models} shows the odds ratio of various effects on a per-user bases. At the strategic tier, doubling the number of drones increases the odds of the user deferring that mission to the agent by 1.6x (p < 0.001), and adding a drone type increases the odds by 1.5x (p < 0.001). At the tactical level, even adding a single drone ... [presumably this too? Maybe?]
\textcolor{red}{We'll remove the last four from associated graph, and check adding just 1 drone. I suspect the tactical might even be significant. If not dw. Also check the above numbers as I read 1.5 and 1.6 off the graph in a hurry. You should have the actual numbers.}

\begin{figure}[H]
\centering
\includegraphics[width=\linewidth,height=0.50\textheight,keepaspectratio]{figs/results/f-size-decides__models__main.png}
\caption{\textbf{Bigger jobs are handed over more, for every size measure and on both tiers.} Each row is a within-person odds ratio per unit of size with its 95\% interval; the bottom block holds the checks for artefacts.}
\label{fig:res:size-decides:models}
\end{figure}

% --- f-size-decides#tactical (pass 5: for the paper) ---

\phantomsection\label{sec:res:size-decides:tactical}\label{sec:res:size-decides:route-people}
Given that missions in this platform progress from strategic to tactical, a key question is where there exists an implied continuity between the two tiers; once the user decides to use the AI assistant at the strategic tier, perhaps that influences their propensity to automate that same mission at the tactical tier. Figure~\ref{fig:res:size-decides:tactical:needed} suggests it does: For the same person and mission size, the odds of doing manual tactical allocation increase by X (p = y). The relative gap is larger for smaller missions, but as the number of drones required increases, this gap narrows. We speculate that this is because the already-established propensity to automate larger tasks becomes a dominating factor and squashes this difference. Participant interviews support this, with one user describing the continuity directly: ``Now I'm using the [tactical] suggestion because I took the [strategic agent's] aggressive approach which means I know that I have enough drones'' (P10). The domination of size as a consideration in when to automate is hinted at also: ``If I look at a mission with 15 drones and I need to move two or three chains, that's just going to take me about a minute, and it's just no point'' (P11).

\begin{figure}[H]
\centering
\includegraphics[width=\linewidth,height=0.50\textheight,keepaspectratio]{figs/results/f-size-decides__tactical__marker.png}
\caption{\textbf{A mission allocated from a strategy card is more likely to get a tactical suggestion.} Share of tactical plans with a suggestion requested, by the drones the mission needs, split by how the mission was allocated.}
\label{fig:res:size-decides:tactical}
\label{fig:res:size-decides:tactical:needed}
\end{figure}

% --- f-size-decides#parallel (pass 5: for the paper) ---
\textcolor{red}{I neglected to mention the strategic. From memory this has similar results. Put the graph back in, and we'll either present both side-by-side, or else focus on one and quote numbers fromthe other. If you can, follow the below paragraph for strat and produce a version for that}
\phantomsection\label{sec:res:size-decides:parallel}
It is worth considering what advantage the agent actually conferred. An analysis only of the qualatative reporting from users would lead one to beleive that the agent is much faster than a human could act. For example [\textcolor{red}{Some quotes here - good ones, and summary of how many mentioned the speed being faster}]. As mentioned in sec [ref back] the agent deliberately takes some time to produce its plan; we aimed for approximate parity with the speed of a well-performing user. Figure~\ref{fig:res:size-decides:parallel} splits each tactical allocation into four phases: Time before action (not shown, but generally slightly higher when by hand), the time taken to build the plan, the time before the plan is approved, and finally the time to actually carry out that plan - i.e. the quality of the solution. The difference in overall time is minimal. Note that the large, by-hand bar only represents n=10(?) events as these tasks were ususally automated, so we advise ignoring this result. The final bar standardises for the relative size differences in which missions were automated by [\textcolor{red}{how?}] and finds a significant speed advantage for using the tactical agent (p = x), but a relatively small effect size ([\textcolor{red}{quantify?}]. We conclude that although many users felt the agent was faster than them, this seems at best to only be a slim advantage. We will discuss in sec \ref{something later} from where this discrepancy may come. \textcolor{red}{Will later mention working in parallel and maybe the hamman time dilation study}

\begin{figure}[H]
\centering
\includegraphics[width=\linewidth,height=0.50\textheight,keepaspectratio]{figs/results/f-size-decides__parallel__all.png}
\caption{\textbf{Hand-built and suggested tactical plans take the operator about the same time; suggested plans are carried out a little sooner.} Mean time per tactical plan in four phases, by mission size.}
\label{fig:res:size-decides:parallel}
\label{fig:res:size-decides:parallel:all}
\end{figure}
\textcolor{red}{This graph should not include the waiting time. The plot below which shoes the forest plot probably won't bake it into the paper, but we will quote a lot of these numbers so make sure they make it in.}

% --- f-size-decides#said (pass 5: for the paper) ---

\phantomsection\label{sec:res:size-decides:said}\label{sec:res:size-decides:parallel-said}
Ten of 33 participants stated the size rule without being asked, and 15 valued the assistants for speed rather than accuracy. Several also described using the wait: while a suggestion filled in, they turned to the next mission. The logs agree: 45\% of tactical plans with a suggestion requested had another mission's strategic panel opened while the planner was open, against 19\% of hand-built plans.
\textcolor{red}{A few lines earlier I said we should lead with the qualitative as this justifies the importance of the graph. This text is the ideal way to do that,}

\clearpage

% ======== f-manual-cards ========

\subsection{Hand-built allocations are anchored on the strategy cards}\label{sec:res:manual-cards}
Something we did not directly plan for was that at the strategic level, operators were asked to build a plan manually, but may do so before or after having seen a suggestion from the agent. This allows us firstly to compare their manual allocations with those of the AI, and secondly to observe whether having seen a suggestion from the agent, this alters the type of manual plan a user may build. Figure~\ref{fig:res:manual-cards:dots} shows every manual assignment relative to the agent's conservative choice. A general trend is seen that they tend to be more extreme: for smaller missions they usually undershoot the AI's choice, and for larger ones they more heavily commit resources. There are a number of confounders to be aware of, for example when few drones were available, some users decided to send the entire reserve: ``All right. Let's just send everything I have. Edit, give me five red ones. Five green ones. Hit it.'' (P10). This on occasion resulted in allocating more than was probably required. \textcolor{red}{I would actually compare this to the "what the hell effect" where people buy more at fast food chains etc but this may be too tenuous a link}

\begin{figure}[H]
\centering
\begin{subfigure}[t]{\linewidth}
\centering
\includegraphics[width=\linewidth,height=0.21\textheight,keepaspectratio]{figs/results/f-manual-cards__dots__all.png}
\caption{All drone types.}
\label{fig:res:manual-cards:dots:all}
\end{subfigure}
\par\smallskip
\end{figure}

% --- f-manual-cards#position (pass 5: for the paper) ---

Although the majority share of manual allocations were under the card suggestions, an effect can be seen when we compare three cases: allocating manually before the AI suggestions appeared (pre-empting), allocating after seeing the cards and rejecting them, or selecting an AI suggestion and editing it (this was done more rarely). Pre-empted allocations were more likely to match what the AI would have done compared to allocations made after seeing the AI suggestions (39\% vs 27\%), and those assignments made after seeing the suggestion were below the AI in 56\% of cases. This may suggest that once a human operator sees an AI suggestion their mind is anchored to it, and they feel more inclined to ``be better than the AI'' [\textcolor{red}{There was a quote about having a grudge against it. Could be good here}]. We should temper this by drawing to the reader's attention that this was not a central focus of the study and our methodology was not entirely robust; we have not been able to reliably determine the degree to which this is an artifact of users intervening \textit{because} the AI over-allocated. More reserach on this question in the use case of drone allocation would be welcome. \textcolor{red}{A link back to some automation bias literature could probably fit well here as it likely has answers to this question - if the literature agrees then we can state our case more strongly.}

\begin{figure}[H]
\centering
\includegraphics[width=\linewidth,height=0.50\textheight,keepaspectratio]{figs/results/f-manual-cards__position__raw.png}
\caption{\textbf{Hand-built allocations mostly sit below both strategy cards or equal to one, rarely above.}}
\label{fig:res:manual-cards:position}
\label{fig:res:manual-cards:position:raw}
\end{figure}

\clearpage

% ======== f-practice ========

\subsection{Learning effects over time}\label{sec:res:practice}
Learning effects are well established in user studies~\cite{TODO}, but are highly relevant in scenarios where a human must team with agents; the speed of improvement in work-rate could inform handover processes in real-life scenarios. We identify both a ``warming up'' effect during a single scenario and a learning effect between the two scenarios (which were separated by a questionnaire and a short break). Figure~\ref{fig:res:practice:speed} shows the median decision time through each session. Decisions sped up within each session, by approximately 6\% per minute at both the strategic and tactical levels. On the strategic tier most of that speed-up was gone by the start of the next session, whereas the tactical speed-up seemed to mostly carry over. It is difficult to conclusively say what caused this discrepancy between the two tiers, although it is possible that increased familiarity and confidence conferred efficiency to the comparatively more straightforward tactical decisions, while their attitudes to the strategic decisions became more complex and considered. \textcolor{red}{May be able to factcheck each of these against the data. Also possibly quotes.}. Participants do note this themselves, with many commenting that they got better as the experiment went on, for example: ``there was definitely an improvement further throughout session one, and in session two, you really find your stride'' (P10). Finally, there is no indication that these differences are related to the relative work-rate throughout the scenario, as shown next.

\begin{figure}[H]
\centering
\begin{subfigure}[t]{\linewidth}
\centering
\includegraphics[width=\linewidth,height=0.21\textheight,keepaspectratio]{figs/results/f-practice__speed__strat.png}
\caption{Strategic tier, sessions side by side.}
\label{fig:res:practice:speed:strat}
\end{subfigure}
\par\smallskip
\begin{subfigure}[t]{\linewidth}
\centering
\includegraphics[width=\linewidth,height=0.21\textheight,keepaspectratio]{figs/results/f-practice__speed__tac.png}
\caption{Tactical tier, sessions side by side.}
\label{fig:res:practice:speed:tac}
\end{subfigure}
\par\smallskip
\caption{\textbf{Faster within a session; on the strategic tier the gain mostly resets by the next.} Median time per decision by pair of minutes: (a) strategic tier, (b) tactical tier.}
\label{fig:res:practice:speed}
\end{figure}

\clearpage

% ======== f-workload ========

\subsection{A busier moment does not push operators onto the assistants}\label{sec:res:workload}
One may naturally assume that a significant factor effecting the operator's propensity to use the AI assistants is the amount of work that is piling up --- that as the mission queue becomes longer, the operator defers more to the agent. This is in fact something which 9/33 participants actively said they did, for example ``If, at first glance, I felt really overwhelmed by the amount of drones to be assigned. I would use the assistant'' (P18), however we did not find any evidence of this, either in them or in the group at large. Figure~\ref{fig:res:workload:raw-vs-adjusted} shows the share of decisions handed over to the assistants by the number of other missions waiting. The dotted lines are the true raw values for the rate of assistant use as the number of other waiting missions increases. The solid line is adjusted for \textcolor{red}{[x, by doing y]}. 


Missions were also handed over no more often after being left waiting in the queue for a long time. Mission that had spend more than two-minutes pending (and hence accumulating a large and visible penalty) were automated 68\% of the time vs 63\% for those that had only just appeared (adjusted for x y z). Tactical suggestions, if anything, became a little less likely the longer a plan had been left open [numbers].


We see no obvious change in the use of the AI assistant as the queue gets longer. This is a perplexing finding and perhaps speaks to a dissonance between perceived and actual action on part of the operator. We will argue later why this may motivate agents which actively help the user manage their use of the assistants. \textcolor{red}{Make sure we do this - it's actually a really nice way to lead into the 'reordering the queue' stuff}

\begin{figure}[H]
\centering
\begin{subfigure}[t]{\linewidth}
\centering
\includegraphics[width=\linewidth,height=0.21\textheight,keepaspectratio]{figs/results/f-workload__raw-vs-adjusted__both-adj.png}
\caption{Both tiers together, adjusted.}
\label{fig:res:workload:raw-vs-adjusted:both-adj}
\end{subfigure}
\par\smallskip
\caption{\textbf{A busier moment does not push operators onto the assistants.} Share of decisions handed over by the number of other missions waiting, both tiers together: (a) raw, with the rate the decisions themselves predict (dashed); (b) adjusted.}
\label{fig:res:workload:raw-vs-adjusted}
\end{figure}
\textcolor{red}{Use a version that shows strat and tac, raw in dotted line and adjusted in solid line.}

\clearpage

% ======== f-flaw-odd ========

\subsection{The illogical strategy card, and who stopped choosing the cards}\label{sec:res:flaw-odd}
Although we did not inject faults into our agents in the truest sense, the simple heuristic design of the strategic agent did mean that it occasionally made recommendations which, while not erroneous, were ``illogical''. Because the conservative policy works by assigning the minimum number of drones the solve the task, and than adds one of each colour (a reasonable policy which might hand-coded by a developer without proper testing), it sometimes made recommendations like the one in figure~\ref{fig:TODO}. These ``illogical cards'' make up 22\% of all recommendations (almost entirely on the simpler, one-type missions), and was a common occurrence for every user, typically very early-on in their first scenario. Conveniently, almost exactly half (15/33) of our users noticed this error and mentioned it, often with confusion: ``Oh, this just sends a bunch of extra of these even though it doesn't need it?'' (P19), and sometimes concluding doubt: ``But in conservative, you don't need camera or lifters, yet it adds it as reserve. Even for tasks it doesn't do. That doesn't make sense.'' (P13). \textcolor{red}{There may be a better second quote here.} This provides us with a useful comparison between those who noticed the error (``noticers'') and those who didn't (``non-noticers), and how it effected their use of that agent.

%\begin{figure}[H]
%\centering
%\includegraphics[TODO}
%\caption{Figure here that shows a bad card suggestion. One with all 3 colours when not needed.}
%\label{fig:TODO}
%\end{figure}

Figure~\ref{fig:res:flaw-odd:event:pairs} compares the noticers and non-noticers' use of the strategic agent over indexed decision (before the first illogical card is shown - the earliest time they could see an issue), that card itself, and future decisions. The illogical card is taken less by both groups. This may indicate the non-noticers didn't actively feel that the agent was making mistakes, but still didn't feel that its recommendation was sensible. Proceeding this, however, there is a marked difference in AI use between the two groups: the noticers' use of the assistant remains X\% lower (significant?) throughout the middle of the scenario, before apparently returning to baseline; by the end of the run, the noticers are using the AI just as much as those who seem to think it is yet to make a mistake. This suggests that there may be some level of trust repair~\cite{TODO}, where having understood that the agent isn't perfect, they become accustomed to it. \textcolor{red}{There may be a quote to help this. If not I think it's ``we speculate that...''}

\begin{figure}[H]
\centering
\includegraphics[width=\linewidth,height=0.50\textheight,keepaspectratio]{figs/results/f-flaw-odd__event__pairs.png}
\caption{\textbf{After the first illogical strategy card, those who later said they noticed it choose the cards less for about ten decisions.} Share of decisions where a strategy card was chosen, by decision number after the first illogical card, in pairs.}
\label{fig:res:flaw-odd:event}
\label{fig:res:flaw-odd:event:pairs}
\end{figure}

After the first scenario, participants answered a few questions about their trust in each agent, and many offered up opinion at this point to the researcher in the room unprompted. We suspect that the process of reflection may have caused them to reconsider their AI use and effected their second scenario. Figure~\ref{fig:res:flaw-odd:by-index:groups} concatenates the two scenarios (meaning whichever one they had first, followed by whichever one they had second: each ``scenario'' is half Strategic-heavy and half Tactical-heavy, grouped together; we are here considering only the ordering) and marks the break. The session 1 data is essentially the same as in Figure~\ref{fig:res:flaw-odd:event:pairs}, but when session 2 begins we see a new pattern: the noticers have once again reduced their usage of the agents, and this remains the case until the end. We suspect this to be the outcome of them having had time to consider, and that they have entered the second scenario with greater scepticism. \textcolor{red}{What about their surveys? Did noticers vs non-noticers have different increases and decreases in trust in the strat agent after each? We could even place this on the graph at the change between 1 and 2.}

\begin{figure}[H]
\centering
\includegraphics[width=\linewidth,height=0.50\textheight,keepaspectratio]{figs/results/f-flaw-odd__by-index__groups.png}
\caption{\textbf{Counted by decision number, the noticers' dip comes before most of them have met an illogical card.} Share of decisions with a strategy card chosen, session 1 then session 2.}
\label{fig:res:flaw-odd:by-index}
\label{fig:res:flaw-odd:by-index:groups}
\end{figure}
\textcolor{red}{Just worth noting that this isn't the same x as the previous figure?. Jsut making sure these are comparable because a reviewer might say ``one of your dots (the fourth blue) here is about 90\%, but in the earlier one none of the dots topped 75\%''. I think the version of this in pairs might be better, but let's have a chat about this.}

% --- f-flaw-odd#spoken (pass 5: for the paper) ---

\clearpage

% ======== f-tier-trust ========

\subsection{The Tactical Assistant is trusted more}\label{sec:res:tier-trust}
Overall, the tactical agent was not only used more, but also clearly more trusted. Only 5 users felt the strategic agent was significantly more trustworthy. Figure~\ref{fig:res:tier-trust:gap} shows the overall trust gap [\textcolor{red}{was this to ``trust in automation scale?'' - if it has a name, use it}] for each user. The trust gap is the effect is moderate (mean $+0.83$ points, $d_z = .62$) [\textcolor{red}{I'm not sure what d\_z means}] and highly significant ($p = .001$). This holds when comparing each of the six questionnaire items for tactical and strategic on a 1-7 scale, i.e. the tactical is rated higher than the stratgic in terms of  ``reliablity'' (+0.94, p<x), ``trust in its recommmendations'' (+1.01, p=y), .... [\textcolor{red}{Here say what each question asked and give the differences plus ideally a p value if that's achievable}]). There were a number of stated reasons for this, but among them were alignment with the human's own strategy: ``The tactical one seemed to be doing what I was doing. So quite reliably. And obviously the strategic one wasn't.'' (P32), and accountability: ``You can't blame the tactical for drone failures but you can blame the strategic for not having redundancy'' (P03).

\begin{figure}[H]
\centering
\includegraphics[width=\linewidth,height=0.50\textheight,keepaspectratio]{figs/results/f-tier-trust__gap__main.png}
\caption{\textbf{The Tactical Assistant is trusted more.} Each person's mean trust in each assistant over their two sessions (1--7).}
\label{fig:res:tier-trust:gap}
\end{figure}

\textcolor{red}{We should also add to the next paragraph a similar comparison of score. Not sure we focus on that so it might as well go here.}
The gap between trust in the two agents was was larger in Strategic Heavy sessions (x for strategic agent vs y for tactical agent; +1.13, p=x) than in Tactical Heavy ones (x for strategic agent vs y for tactical agent; +0.52, p=x)


\clearpage
